\section{Network and Threat Model}
\label{sec:model}

\subsection{Network model}

The network is a dynamic graph $G(t) = (V, E(t))$ over $N$ mobile nodes with
time-varying wireless links. A subset $\mathcal{B} \subseteq V$ consists of beacon
nodes advertising position information usable by neighbours for
localisation. A subset $\mathcal{B}_m \subseteq \mathcal{B}$ is compromised.
Each monitored beacon $i$ carries a state
\begin{equation}
\label{eq:state}
S_i(t) \in \{\textsc{Normal}, \textsc{Suspicious}, \textsc{Confirmed}\}.
\end{equation}
The intermediate state exists so that a single anomalous observation does
not directly authorise removal.

The pipeline both raises suspicion and decides what to do about it; our
contribution concerns the second step, verification and response. Attackers are insiders
holding valid credentials, so cryptographic authentication does not resolve
the problem: an authenticated node can still lie.

\subsection{Attack model}

We evaluate three attacker behaviours, chosen so the first two are disjoint
in the evidence they disturb and the third combines them.

\textbf{A1, false location.} The compromised beacon advertises a position
displaced from its true position, while forwarding all transit traffic
honestly. This corrupts localisation and, by construction, produces no
forwarding anomaly. The offset $\delta$ is applied to both coordinates and
the result is clamped to the field,
\begin{equation}
\label{eq:clamp}
x^{c} = \min\bigl(X, \max(0, x + \delta)\bigr), \qquad
y^{c} = \min\bigl(Y, \max(0, y + \delta)\bigr).
\end{equation}
An unclamped beacon is therefore displaced by $\delta\sqrt{2} = 424$\,m for
$\delta = 300$\,m, while one near the far corner is displaced much less.
Over the campaign the realised displacement averaged 267\,m, with median
271\,m and 95th percentile 299\,m (Section~\ref{sec:locerr}). We report the
measured displacement alongside the nominal parameter, since the effective
attack strength depends on where in the field the attacker happens to be.

\textbf{A2, selective forwarding.} The compromised beacon advertises its true
position and drops transit packets with probability $p$. This degrades
delivery and produces no localisation anomaly.

\textbf{A3, hybrid at half intensity.} The compromised beacon applies both
behaviours with each halved: offset $\delta/2$ and drop probability $p/2$.

The choice to halve rather than compound is deliberate. Compounding both
behaviours at full strength produces an attacker that is more conspicuous on
every axis at once, which tests nothing A1 and A2 do not already test.
Halving produces an attacker that is less conspicuous on each axis while
remaining malicious on both. That is the configuration under which a fusion
advantage, if the mechanism confers one, should appear, because each
single-evidence branch operates closer to its detection threshold. A3 is
therefore a test of the fusion hypothesis and not a demonstration of it,
and Section~\ref{sec:sweep} sweeps its intensity.

\begin{table}[!t]
\caption{Principal notation.}
\label{tab:notation}
\centering
\small
\begin{tabular}{@{}ll@{}}
\toprule
Symbol & Meaning \\
\midrule
$FR_i(t)$ & forwarding ratio of beacon $i$ \\
$z_B, z_L$ & standardised behavioural / localisation evidence \\
$A_B, A_L$ & oriented anomaly scores \\
$Z_{ij}(w)$ & fused evidence statistic \\
$T^D, T^I, T^H$ & direct, indirect, historical trust \\
$\hat C_i(t)$ & collapsed confidence \\
$\tau_N, \tau_M$ & normal and malicious thresholds \\
$W, K$ & persistence window and confirmation count \\
$\mathcal{B}_{iso}(t)$ & isolation set \\
$\theta_B, \theta_L$ & branch detection thresholds in attack intensity \\
$R_{TH}$ & throughput recovery ratio \\
$\Phi, \bar\Phi$ & cumulative false isolation (node-seconds) and its normalised form \\
\bottomrule
\end{tabular}
\end{table}
